\documentclass[11pt]{amsart}
\usepackage{geometry}                % See geometry.pdf to learn the layout options. There are lots.
\geometry{letterpaper}                   % ... or a4paper or a5paper or ... 
%\geometry{landscape}                % Activate for for rotated page geometry
%\usepackage[parfill]{parskip}    % Activate to begin paragraphs with an empty line rather than an indent
\usepackage{graphicx}
\usepackage{amssymb}
\usepackage{epstopdf}
\DeclareGraphicsRule{.tif}{png}{.png}{`convert #1 `dirname #1`/`basename #1 .tif`.png}

\title{Ejemplo 3.3}
%\date{}                                           % Activate to display a given date or no date

\begin{document}
\maketitle
%\section{}
%\subsection{}
\noindent Escriba un diagrama de flujo tal que dado como datos N números enteros, ob­ tenga el número de ceros que hay entre estos números. \newline
 \noindent  \[Datos: N,
\text{ NUMi}: NUM_1, NUM_2, \ldots, NUM_n
\]
 \noindent  Donde:  \newline
 
N es una variable de tipo entero que representa el número de datos que se ingresan.\newline
\[
\text{NUMj}: \text{{es una variable de tipo entera que representa al número }} i \quad (1 \leq i \leq N)
\]
\noindent  \textbf{Explicación de las variables} \newline

 \noindent  I: Variable de tipo entero. Representa al contador del ciclo. \newline
N y NUM: Variables de tipo entero.\newline
CUECER: Variable de tipo entero. Cuenta el número de ceros.\newline

A continuación en la siguiente tabla podemos observar el seguimiento del al­goritmo.\newline
\begin{table}[htbp]
\centering
\label{tab:corrida}
\begin{tabular}{|c|c|c|c|c|c|}
\hline
\multicolumn{4}{ |c| }{\textbf {Tabla 3.2}} \\ \hline
\textbf{l} & \textbf{M} & \textbf{CUECER} & \textbf{NUM}  \\ \hline
1 & 12  & 0 & \\ \hline
2 &   &  & 18 \\ \hline
3 &   &  & 23\\ \hline
4 &   &  1 & 0 \\ \hline
5 &   &  & 17\\ \hline
6 &   &  & 22 \\ \hline
7&   &  2& 0 \\ \hline
8&   &  & 37 \\ \hline
9&   &  & 43 \\ \hline
10&   & 3 & 0 \\ \hline
11&   &  & 27 \\ \hline
12&   &  & 41 \\ \hline
13 &   &  & 53 \\ \hline
\end{tabular}
\end{table}
\end{document}  